% GENERATED FILE. Edit canonical lessons, then run npm run generate:books.
% canonical-source-hash: fnv1a64:a8d87230a577bbd7
% canonical-lessons: MR-C477-talhat, MR-C477-muth, MR-C477-mandi, MR-C477-potri, MR-C477-ghota

\chapter{Hands and Legs}
\label{ch:mr-hands-and-legs}
\hlchaptermodality{477}

\begin{chapteropening}
\textbf{By the end of this chapter:} \emph{I can name the palm, a fist, the thigh, the calf and the ankle.}

Complete the last lesson of chapter 477: I can name the palm, a fist, the thigh, the calf and the ankle.
\end{chapteropening}

\section[\texorpdfstring{\mr{तळहात} (taḷhāt)}{तळहात (taḷhāt)}]{\mr{तळहात} (taḷhāt) — the palm}

\label{lesson:MR-C477-talhat}



\begin{quote}
Before the new one: say the Marathi for a beard, then the Marathi for a moustache.
\end{quote}

\subsection*{You'll want to know: \mr{तळहात}}
\textbf{\mr{तळहात}} — \emph{taḷhāt} — \textquotedblleft{}the palm\textquotedblright{}.

\begin{quote}
\mr{आजीने} \mr{माझ्या} \mr{तळहातावर} \mr{साखर} \mr{ठेवली}.
\end{quote}

\begin{quote}
\emph{ājīne mājhyā taḷhātāvar sākhar ṭhevlī.}
\end{quote}

\begin{quote}
Grandmother put some sugar on my palm.
\end{quote}

Masculine.

\mr{तळ}, the base or flat, + \mr{हात}. The sole of the foot is \mr{तळपाय}.

\begin{cousinweb}
\textbf{\mr{तळहात}} (\emph{taḷhāt}) is a Marathi compound of \mr{तळ}, \textquotedblleft{}base, flat\textquotedblright{}, and \mr{हात}, \textquotedblleft{}hand\textquotedblright{}, a tadbhava of Sanskrit \emph{hasta} (CDIAL 14024).
\end{cousinweb}

\begin{grammarlens}[title={What you've built}]
One more word for describing what you lived and felt.
\end{grammarlens}

\subsection*{Your turn}
\begin{itemize}
\raggedright
  \item \emph{Say it:} \emph{taḷhāt}
  \item \emph{Say it:} \emph{taḷhāt}, once more
  \item \emph{Say it:} \emph{ājīne mājhyā taḷhātāvar sākhar ṭhevlī.} — Grandmother put some sugar on my palm.
  \item \emph{Recall:} say the Marathi for a beard, then the Marathi for a moustache, then say \emph{taḷhāt} again
\end{itemize}

\begin{tcolorbox}[breakable,colback=teal!4,colframe=teal!35!black,title={Before you move on}]
What does \mr{तळहात} mean? (\textquotedblleft{}The palm\textquotedblright{}.) Say it once more.
\end{tcolorbox}
\section[\texorpdfstring{\mr{मूठ} (mūṭh)}{मूठ (mūṭh)}]{\mr{मूठ} (mūṭh) — a fist; a handful}

\label{lesson:MR-C477-muth}



\begin{quote}
Before the new one: say the Marathi for a moustache, then the Marathi for the palm.
\end{quote}

\subsection*{You'll want to know: \mr{मूठ}}
\textbf{\mr{मूठ}} — \emph{mūṭh} — \textquotedblleft{}a fist; a handful\textquotedblright{}.

\begin{quote}
\mr{त्याने} \mr{मुठीत} \mr{एक} \mr{नाणं} \mr{लपवलं}.
\end{quote}

\begin{quote}
\emph{tyāne muṭhīt ek nāṇaṁ lapavlaṁ.}
\end{quote}

\begin{quote}
He hid a coin in his fist.
\end{quote}

Feminine.

Oblique \mr{मुठी}-: \mr{मुठीत}, in the fist. \mr{मूठभर}, a handful.

\begin{cousinweb}
\textbf{\mr{मूठ}} (\emph{mūṭh}) is a tadbhava, inherited from Sanskrit \emph{muṣṭi}, \textquotedblleft{}fist\textquotedblright{}; CDIAL lists Marathi \emph{mūṭh} (CDIAL 10221).
\end{cousinweb}

\begin{grammarlens}[title={What you've built}]
One more word for describing what you lived and felt.
\end{grammarlens}

\subsection*{Your turn}
\begin{itemize}
\raggedright
  \item \emph{Say it:} \emph{mūṭh}
  \item \emph{Say it:} \emph{mūṭh}, once more
  \item \emph{Say it:} \emph{tyāne muṭhīt ek nāṇaṁ lapavlaṁ.} — He hid a coin in his fist.
  \item \emph{Recall:} say the Marathi for a moustache, then the Marathi for the palm, then say \emph{mūṭh} again
\end{itemize}

\begin{tcolorbox}[breakable,colback=teal!4,colframe=teal!35!black,title={Before you move on}]
What does \mr{मूठ} mean? (\textquotedblleft{}A fist; a handful\textquotedblright{}.) Say it once more.
\end{tcolorbox}
\section[\texorpdfstring{\mr{मांडी} (māṇḍī)}{मांडी (māṇḍī)}]{\mr{मांडी} (māṇḍī) — the thigh; the lap}

\label{lesson:MR-C477-mandi}



\begin{quote}
Before the new one: say the Marathi for the palm, then the Marathi for a fist.
\end{quote}

\subsection*{You'll want to know: \mr{मांडी}}
\textbf{\mr{मांडी}} — \emph{māṇḍī} — \textquotedblleft{}the thigh; the lap\textquotedblright{}.

\begin{quote}
\mr{बाळ} \mr{आईच्या} \mr{मांडीवर} \mr{झोपलं}.
\end{quote}

\begin{quote}
\emph{bāḷ āīcyā māṇḍīvar jhoplaṁ.}
\end{quote}

\begin{quote}
The baby fell asleep on its mother's lap.
\end{quote}

Feminine.

\mr{मांडीवर}, on the lap; \mr{मांडी} \mr{घालून} \mr{बसणे}, to sit cross-legged.

\begin{cousinweb}
\textbf{\mr{मांडी}} (\emph{māṇḍī}) is a Marathi noun whose origin is not settled.
\end{cousinweb}

\begin{grammarlens}[title={What you've built}]
One more word for describing what you lived and felt.
\end{grammarlens}

\subsection*{Your turn}
\begin{itemize}
\raggedright
  \item \emph{Say it:} \emph{māṇḍī}
  \item \emph{Say it:} \emph{māṇḍī}, once more
  \item \emph{Say it:} \emph{bāḷ āīcyā māṇḍīvar jhoplaṁ.} — The baby fell asleep on its mother's lap.
  \item \emph{Recall:} say the Marathi for the palm, then the Marathi for a fist, then say \emph{māṇḍī} again
\end{itemize}

\begin{tcolorbox}[breakable,colback=teal!4,colframe=teal!35!black,title={Before you move on}]
What does \mr{मांडी} mean? (\textquotedblleft{}The thigh; the lap\textquotedblright{}.) Say it once more.
\end{tcolorbox}
\section[\texorpdfstring{\mr{पोटरी} (poṭrī)}{पोटरी (poṭrī)}]{\mr{पोटरी} (poṭrī) — the calf (of the leg)}

\label{lesson:MR-C477-potri}



\begin{quote}
Before the new one: say the Marathi for a fist, then the Marathi for the thigh.
\end{quote}

\subsection*{You'll want to know: \mr{पोटरी}}
\textbf{\mr{पोटरी}} — \emph{poṭrī} — \textquotedblleft{}the calf (of the leg)\textquotedblright{}.

\begin{quote}
\mr{डास} \mr{माझ्या} \mr{पोटरीला} \mr{चावला}.
\end{quote}

\begin{quote}
\emph{ḍās mājhyā poṭrīlā cāvlā.}
\end{quote}

\begin{quote}
A mosquito bit me on the calf.
\end{quote}

Feminine.

Feminine: \mr{पोटरीला}, on the calf.

\begin{cousinweb}
\textbf{\mr{पोटरी}} (\emph{poṭrī}) is a Marathi noun whose origin is not settled.
\end{cousinweb}

\begin{grammarlens}[title={What you've built}]
One more word for describing what you lived and felt.
\end{grammarlens}

\subsection*{Your turn}
\begin{itemize}
\raggedright
  \item \emph{Say it:} \emph{poṭrī}
  \item \emph{Say it:} \emph{poṭrī}, once more
  \item \emph{Say it:} \emph{ḍās mājhyā poṭrīlā cāvlā.} — A mosquito bit me on the calf.
  \item \emph{Recall:} say the Marathi for a fist, then the Marathi for the thigh, then say \emph{poṭrī} again
\end{itemize}

\begin{tcolorbox}[breakable,colback=teal!4,colframe=teal!35!black,title={Before you move on}]
What does \mr{पोटरी} mean? (\textquotedblleft{}The calf (of the leg)\textquotedblright{}.) Say it once more.
\end{tcolorbox}
\section[\texorpdfstring{\mr{घोटा} (ghoṭā)}{घोटा (ghoṭā)}]{\mr{घोटा} (ghoṭā) — the ankle}

\label{lesson:MR-C477-ghota}



\begin{quote}
Before the new one: say the Marathi for the thigh, then the Marathi for the calf.
\end{quote}

\subsection*{You'll want to know: \mr{घोटा}}
\textbf{\mr{घोटा}} — \emph{ghoṭā} — \textquotedblleft{}the ankle\textquotedblright{}.

\begin{quote}
\mr{जिन्यावरून} \mr{उतरताना} \mr{माझा} \mr{घोटा} \mr{वळला}.
\end{quote}

\begin{quote}
\emph{jinyāvarūn utartānā mājhā ghoṭā vaḷlā.}
\end{quote}

\begin{quote}
My ankle turned as I came down the stairs.
\end{quote}

Masculine.

Masculine: \mr{घोट्याला}, at the ankle.

\begin{cousinweb}
\textbf{\mr{घोटा}} (\emph{ghoṭā}) is a Marathi noun whose origin is not settled.
\end{cousinweb}

\begin{grammarlens}[title={What you've built}]
One more word for describing what you lived and felt.
\end{grammarlens}

\subsection*{Your turn}
\begin{itemize}
\raggedright
  \item \emph{Say it:} \emph{ghoṭā}
  \item \emph{Say it:} \emph{ghoṭā}, once more
  \item \emph{Say it:} \emph{jinyāvarūn utartānā mājhā ghoṭā vaḷlā.} — My ankle turned as I came down the stairs.
  \item \emph{Recall:} say the Marathi for the thigh, then the Marathi for the calf, then say \emph{ghoṭā} again
\end{itemize}

\begin{tcolorbox}[breakable,colback=teal!4,colframe=teal!35!black,title={Before you move on}]
What does \mr{घोटा} mean? (\textquotedblleft{}The ankle\textquotedblright{}.) Say it once more.
\end{tcolorbox}
